% Part: lambda-calculus
% Chapter: church-rosser
% Section: beta-reduction

\documentclass[../../../include/open-logic-section]{subfiles}

\begin{document}

\olfileid{lam}{cr}{b}

\olsection{اختزال~$\beta$}

\begin{lem}\ollabel{lem:one-par}
  من~${\OLId{latin-looped}{M}} \bredone {\OLId{latin-looped}{M}}'$ يلزم~${\OLId{latin-looped}{M}} \bredpar {\OLId{latin-looped}{M}}'$.
\end{lem}
\begin{proof}
  نستقرئ اشتقاق~${\OLId{latin-looped}{M}} \bredone {\OLId{latin-looped}{M}}'$ من أصغر علاقة متوافقة، فتقع أربع حالات:
  \begin{enumerate}
  \item في تقلص الجذر يكون~${\OLId{latin-looped}{M}}=(\lambd[{\OLId{latin-ordinary}{x}}][{\OLId{latin-looped}{N}}]){\OLId{latin-looped}{Q}}$ و
    ${\OLId{latin-looped}{M}}'=\Subst{{\OLId{latin-looped}{N}}}{{\OLId{latin-looped}{Q}}}{{\OLId{latin-ordinary}{x}}}$. ولما كان~${\OLId{latin-looped}{N}}\bredpar {\OLId{latin-looped}{N}}$ و${\OLId{latin-looped}{Q}}\bredpar {\OLId{latin-looped}{Q}}$
    بموجب \olref[pb]{thm:refl}، أعطت
    \olref[pb]{defn:bredpar}\olref[pb]{defn:bredpar4} المطلوب.
  \item وإن وقع التقلص تحت تجريد، كان~${\OLId{latin-looped}{M}}=\lambd[{\OLId{latin-ordinary}{x}}][{\OLId{latin-looped}{P}}]$ و
    ${\OLId{latin-looped}{M}}'=\lambd[{\OLId{latin-ordinary}{x}}][{\OLId{latin-looped}{P}}']$ مع~${\OLId{latin-looped}{P}}\bredone {\OLId{latin-looped}{P}}'$. يفيد فرض الاستقراء
    ${\OLId{latin-looped}{P}}\bredpar {\OLId{latin-looped}{P}}'$، ثم تفيد \olref[pb]{defn:bredpar2} المطلوب.
  \item وإن وقع في الطرف الأيسر من تطبيق~$PQ$، جمعنا فرض
    الاستقراء لذلك الطرف إلى~${\OLId{latin-looped}{Q}}\bredpar {\OLId{latin-looped}{Q}}$ بالانعكاس، ثم طبقنا
    \olref[pb]{defn:bredpar3}.
  \item وإن وقع في الطرف الأيمن، جمعنا~${\OLId{latin-looped}{P}}\bredpar {\OLId{latin-looped}{P}}$ بالانعكاس
    إلى فرض الاستقراء للطرف الأيمن، ثم طبقنا القاعدة نفسها.
  \end{enumerate}
  % تصحيح مصدر (0370-I1): شمل الاستقراء حالات التوافق الثلاث فضلًا على تقلص الجذر.
\end{proof}

\begin{lem}\ollabel{lem:par-red}
  من~${\OLId{latin-looped}{M}} \bredpar {\OLId{latin-looped}{M}}'$ يلزم~${\OLId{latin-looped}{M}} \bred {\OLId{latin-looped}{M}}'$.
\end{lem}

\begin{proof}
  نستقرئ~!!{derivation} العبارة~${\OLId{latin-looped}{M}} \bredpar {\OLId{latin-looped}{M}}'$، ونفصل بحسب آخر قاعدة:
  \begin{enumerate}
  \item إن كانت~\olref[pb]{defn:bredpar1}، لم يكن~${\OLId{latin-looped}{M}}$ و${\OLId{latin-looped}{M}}'$ إلا~${\OLId{latin-ordinary}{x}}$،
    ولدينا~${\OLId{latin-ordinary}{x}} \bred {\OLId{latin-ordinary}{x}}$.
  \item وإن كانت~\olref[pb]{defn:bredpar2}، كان~${\OLId{latin-looped}{M}}$ هو~$\lambd[{\OLId{latin-ordinary}{x}}][{\OLId{latin-looped}{N}}]$
    و${\OLId{latin-looped}{M}}'$ هو~$\lambd[{\OLId{latin-ordinary}{x}}][{\OLId{latin-looped}{N}}']$، لبعض~${\OLId{latin-ordinary}{x}}$ و${\OLId{latin-looped}{N}}$ و${\OLId{latin-looped}{N}}'$، مع~${\OLId{latin-looped}{N}} \bredpar {\OLId{latin-looped}{N}}'$.
    ومن فرض الاستقراء~${\OLId{latin-looped}{N}} \bred {\OLId{latin-looped}{N}}'$؛ فتنتج~$\lambd[{\OLId{latin-ordinary}{x}}][{\OLId{latin-looped}{N}}] \bred \lambd[{\OLId{latin-ordinary}{x}}][{\OLId{latin-looped}{N}}']$
    بإجراء تقلصات~$\bredone$ نفسها التي أعطت~${\OLId{latin-looped}{N}} \bred {\OLId{latin-looped}{N}}'$.
  \item وإن كانت~\olref[pb]{defn:bredpar3}، كان~${\OLId{latin-looped}{M}}$ هو~$PQ$ و${\OLId{latin-looped}{M}}'$ هو~${\OLId{latin-looped}{P}}'{\OLId{latin-looped}{Q}}'$،
    لبعض~${\OLId{latin-looped}{P}}$ و${\OLId{latin-looped}{Q}}$ و${\OLId{latin-looped}{P}}'$ و${\OLId{latin-looped}{Q}}'$، مع~${\OLId{latin-looped}{P}} \bredpar {\OLId{latin-looped}{P}}'$ و${\OLId{latin-looped}{Q}} \bredpar {\OLId{latin-looped}{Q}}'$.
    ويفيد فرض الاستقراء~${\OLId{latin-looped}{P}} \bred {\OLId{latin-looped}{P}}'$ و${\OLId{latin-looped}{Q}} \bred {\OLId{latin-looped}{Q}}'$.
    فنحصل على~$PQ \bred {\OLId{latin-looped}{P}}'{\OLId{latin-looped}{Q}}'$ بإجراء الاختزال~${\OLId{latin-looped}{P}} \bred {\OLId{latin-looped}{P}}'$
    ثم الاختزال~${\OLId{latin-looped}{Q}} \bred {\OLId{latin-looped}{Q}}'$.
  \item وإن كانت~\olref[pb]{defn:bredpar4}، كان~${\OLId{latin-looped}{M}}$ هو~$(\lambd[{\OLId{latin-ordinary}{x}}][{\OLId{latin-looped}{N}}]){\OLId{latin-looped}{Q}}$
    و${\OLId{latin-looped}{M}}'$ هو~$\Subst{{\OLId{latin-looped}{N}}'}{{\OLId{latin-looped}{Q}}'}{{\OLId{latin-ordinary}{x}}}$، لبعض~${\OLId{latin-ordinary}{x}}$ و${\OLId{latin-looped}{N}}$ و${\OLId{latin-looped}{N}}'$ و${\OLId{latin-looped}{Q}}$ و${\OLId{latin-looped}{Q}}'$،
    مع~${\OLId{latin-looped}{N}} \bredpar {\OLId{latin-looped}{N}}'$ و${\OLId{latin-looped}{Q}} \bredpar {\OLId{latin-looped}{Q}}'$.
    ومن فرض الاستقراء~${\OLId{latin-looped}{Q}} \bred {\OLId{latin-looped}{Q}}'$ و${\OLId{latin-looped}{N}} \bred {\OLId{latin-looped}{N}}'$.
    فبلوغ~$(\lambd[{\OLId{latin-ordinary}{x}}][{\OLId{latin-looped}{N}}]){\OLId{latin-looped}{Q}} \bred \Subst{{\OLId{latin-looped}{N}}'}{{\OLId{latin-looped}{Q}}'}{{\OLId{latin-ordinary}{x}}}$
    يكون بإجراء~${\OLId{latin-looped}{N}} \bred {\OLId{latin-looped}{N}}'$، ثم~${\OLId{latin-looped}{Q}} \bred {\OLId{latin-looped}{Q}}'$،
    ثم تقليص~$(\lambd[{\OLId{latin-ordinary}{x}}][{\OLId{latin-looped}{N}}']){\OLId{latin-looped}{Q}}'$ إلى~$\Subst{{\OLId{latin-looped}{N}}'}{{\OLId{latin-looped}{Q}}'}{{\OLId{latin-ordinary}{x}}}$.
  \end{enumerate}
\end{proof}

\begin{lem}\ollabel{lem:str}
  العلاقة~$\bred$ أصغر علاقة متعدية تشتمل على~$\bredpar$.
\end{lem}

\begin{proof}
  لتكن~$\xred$ أصغر العلاقات المتعدية التي تشتمل على~$\bredpar$.

  أما~$\bred \subseteq \xred$، فلنفرض~${\OLId{latin-looped}{M}} \bred {\OLId{latin-looped}{M}}'$، أي
  ${\OLId{latin-looped}{M}} \ident {\OLId{latin-looped}{M}}_{\OLMathNumeral{1}} \bredone \dots \bredone {\OLId{latin-looped}{M}}_{\OLId{latin-ordinary}{k}} \ident {\OLId{latin-looped}{M}}'$.
  وتفيد~\olref{lem:one-par} أن~${\OLId{latin-looped}{M}} \ident {\OLId{latin-looped}{M}}_{\OLMathNumeral{1}} \bredpar \dots \bredpar {\OLId{latin-looped}{M}}_{\OLId{latin-ordinary}{k}} \ident {\OLId{latin-looped}{M}}'$.
  ولما كانت~$\xred$ تشتمل على~$\bredpar$ وتتصف بالتعدي، نتج~${\OLId{latin-looped}{M}} \xred {\OLId{latin-looped}{M}}'$.

  وأما~$\xred \subseteq \bred$، فلنفرض~${\OLId{latin-looped}{M}} \xred {\OLId{latin-looped}{M}}'$، أي
  ${\OLId{latin-looped}{M}} \ident {\OLId{latin-looped}{M}}_{\OLMathNumeral{1}} \bredpar \dots \bredpar {\OLId{latin-looped}{M}}_{\OLId{latin-ordinary}{k}} \ident {\OLId{latin-looped}{M}}'$.
  وتفيد~\olref{lem:par-red} أن~${\OLId{latin-looped}{M}} \ident {\OLId{latin-looped}{M}}_{\OLMathNumeral{1}} \bred \dots \bred {\OLId{latin-looped}{M}}_{\OLId{latin-ordinary}{k}} \ident {\OLId{latin-looped}{M}}'$.
  ومن تعدي~$\bred$ ينتج~${\OLId{latin-looped}{M}} \bred {\OLId{latin-looped}{M}}'$.
\end{proof}

\begin{thm}\ollabel{thm:cr}
  العلاقة~$\bred$ ذات خاصية تشيرش--روسر.
\end{thm}

\begin{proof}
  تجمع~\olref[dap]{thm:str} و\olref[pb]{thm:cr} و\olref{lem:str}
  فتنتج المبرهنة مباشرة.
\end{proof}

\end{document}
